\Needspace{16\baselineskip}
\fancyhead[L]{\small GAUSS AN VON LINDENAU.}
\begin{center}
GAUSS AN VON LINDENAU.
\end{center}

\hrule
\begin{center}
\fontsize{10}{12}\selectfont Monatliche Correspondenz. Band XVIII. S. 269..273. September 1808.
\end{center}
\hrule

\begin{flushright}
Göttingen, am 30. August 1808.
\end{flushright}

\par\indent Ich danke Ihnen sehr herzlich dafür, dass Sie mir die Nachrichten über die beiden letzten Kometen und Ihre Beobachtungen der neuen Planeten so freundlich mitgeteilt haben. Schade, dass die Nachrichten nicht früher bekannt geworden sind. Jetzt ist natürlich an eine Suche nach den Kometen nicht mehr zu denken. Gestatten Sie mir, Ihnen die neuesten Ergebnisse unserer Juno-Beobachtungen im beiliegenden Blatt unserer gel. Anz. vorzulegen. Die Beobachtungen waren folgende:


\begin{center}\fontsize{10}{12}\selectfont\setlength{\tabcolsep}{3pt}\renewcommand{\arraystretch}{1.15}\begin{tabular}{lrrr}
1808 & \shortstack{Mittlere Zeit\\in Göttingen} & \shortstack{Scheinbare\\Rektaszension der Juno} & \shortstack{Südliche Deklination\\der Juno}\\\hline
Juni 20 & $11^{\text{u}}\;49^{\text{m}}\;0^{\text{s}}$ & $315^\circ\;29'\;34''$ & $2^\circ\;16'\;23''$\\
22 & $12\;0\;45$ & $315\;23\;1$ & $2\;14\;29$\\
Juli 6 & $12\;42\;23$ & $314\;1\;0$ & $2\;19\;7$\\
30 & $10\;51\;17$ & $309\;38\;0$ & $3\;57\;13$\\
31 & $11\;40\;57$ & $309\;24\;53$ & $4\;3\;48$\\
Aug. 4 & $10\;42\;53$ & $308\;33\;44$ & $4\;30\;32$\\
5 & $10\;53\;46$ & $308\;20\;39$ & $4\;37\;36$\\
\end{tabular}\end{center}


\Needspace{11\baselineskip}
\par\indent Demnach ist der Fehler der Ephemeride bei den letzten Beobachtungen in der Rektaszension auf $13'$ angewachsen. Der Fehler in der Deklination liegt ziemlich unverändert bei $1'5$. Beobachtungen von anderen Orten sind bisher noch nicht bekannt geworden. Die letzten vier Beobachtungen ergeben für die Opposition folgendes Resultat:

\begin{center}
1808 August 2. $9^{\text{u}}\,30^{\text{m}}\,43^{\text{s}}$ mittlere Zeit in Göttingen\\
Wahre Länge .\;.\;.\;.\;.\;.\;.\;.\;.\; $310^\circ\,16'\,31''5$\\
Wahre geozentrische Breite .\;.\;.\;.\; $13\;53\;56,5$ nordl. \end{center}

\par\indent Diese Opposition, die vierte bisher beobachtete, wurde anschließend mit denen von 1804, 1806 und 1807 verbunden, um die folgenden neuen Elemente (VIII) zu bestimmen. Sie geben sämtliche bisherigen Beobachtungen weiterhin sehr gut wieder, [sie sind oben auf Seite 306 aus den Gött. gel. Anzeigen vom 25. Aug. 1808 abgedruckt.]

\par\indent Vesta haben wir seit dem 12. August schon achtmal im Meridian beobachtet, aber die Beobachtungen sind noch nicht alle berechnet.

\par\indent [Den weiteren Teil des Briefes, der den Aufsatz über die Bestimmung der Polhöhe ankündigt, habe ich oben auf Seite 146 dieses Bandes nach jenem Aufsatz wiedergegeben. \textsc{Schering}.]
